% =========================================================
% Proof Strategies for Quantified Statements
% =========================================================

\section{Proof Strategies for Quantified Statements}

% ---------------------------------------------------------
% TOOLKIT
% ---------------------------------------------------------
\begin{tcolorbox}[colback=gray!6, colframe=gray!40, arc=2pt,
  left=6pt, right=6pt, top=4pt, bottom=4pt,
  title={\small\textbf{Quantifier Proof Strategies — Quick Reference}},
  fonttitle=\small\bfseries]
\small
\begin{tabular}{l l l}
\toprule
\textbf{Goal shape} & \textbf{Strategy} & \textbf{Key rule} \\
\midrule
Prove $\forall x\,\varphi(x)$  & Fix an arbitrary $x$; prove $\varphi(x)$ for that $x$  & UG \\
Prove $\exists x\,\varphi(x)$  & Produce a specific witness $t$; verify $\varphi(t)$     & EG \\
Use $\forall x\,\varphi(x)$    & Instantiate at any term $t$ to get $\varphi(t)$         & UI \\
Use $\exists x\,\varphi(x)$    & Introduce a fresh name $c$ for the witness             & EI \\
Negate $\forall x\,\varphi$    & Becomes $\exists x\,\neg\varphi$                        & Neg.\ law \\
Negate $\exists x\,\varphi$    & Becomes $\forall x\,\neg\varphi$                        & Neg.\ law \\
\midrule
Prove $\forall x\,\exists y\,\varphi(x,y)$ & Fix $x$; then construct $y$ depending on $x$ & UG then EG \\
Prove $\exists y\,\forall x\,\varphi(x,y)$ & Find $y$ first; then verify for all $x$       & EG then UG \\
\bottomrule
\end{tabular}
\end{tcolorbox}

\vspace{1em}

The inference rules UI, UG, EI, and EG were defined in the previous section.
That section answered the question: what does each rule license? This section
answers a different question: when the goal of a proof has a particular
quantifier structure, which rule should be used, and what does the proof
skeleton look like? The two questions are distinct, and both must be understood
clearly before reading analysis proofs.

The single most important organizing principle is this: \emph{the shape of the
goal determines the proof structure, and the shape of the hypothesis determines
how to use it}. A $\forall$ goal calls for a different move than a $\exists$
goal; a $\forall$ hypothesis is used differently from a $\exists$ hypothesis.
The four combinations are the four strategies below.

% ---------------------------------------------------------
% Proving a universal
% ---------------------------------------------------------

\section*{Strategy 1: Proving a Universal $\forall x\,\varphi(x)$}

To prove $\forall x\,\varphi(x)$: introduce a fresh variable $x_0$ representing
an \emph{arbitrary, unspecified} element of the domain. Prove $\varphi(x_0)$
without making any assumption about $x_0$ beyond its being in the domain. Then
apply Universal Generalization (UG) to conclude $\forall x\,\varphi(x)$.

\begin{quote}
\textbf{Let} $x_0$ be an arbitrary element of the domain.\\
\textbf{[prove $\varphi(x_0)$ using only the hypotheses, not any special property of $x_0$]}\\
\textbf{Therefore:} $\forall x\,\varphi(x)$. \quad (by UG)
\end{quote}

The critical constraint is that $x_0$ must be \emph{arbitrary}: no hypothesis
in the proof may impose a condition on $x_0$ that is not universally true.
Violating this constraint is the UG fallacy (``illicit generalization'').

\begin{remark}[Why ``let $x$ be arbitrary'' works]
Fixing an arbitrary element $x_0$ is the formal counterpart of the English
phrase ``let $x$ be any element.'' The word ``arbitrary'' means: the proof
uses no information about $x_0$ that is specific to $x_0$ --- only information
that applies to every element of the domain. If the argument goes through for
such an unspecified $x_0$, it goes through for every $x$, so the universal
follows.
\end{remark}

\begin{remark}[The UG side condition in practice]
UG requires that $x_0$ not appear free in any undischarged assumption. In
practice: do not use any hypothesis of the form ``$x_0$ has property $P$''
unless $P$ holds of every domain element. If a hypothesis says ``$x_0 > 0$,''
then $x_0$ is not arbitrary, and UG to $\forall x$ is invalid.
\end{remark}

% ---------------------------------------------------------
% Proving an existential
% ---------------------------------------------------------

\section*{Strategy 2: Proving an Existential $\exists x\,\varphi(x)$}

To prove $\exists x\,\varphi(x)$: exhibit a specific term $t$ and verify that
$\varphi(t)$ holds. Then apply Existential Generalization (EG) to conclude
$\exists x\,\varphi(x)$.

\begin{quote}
\textbf{Claim:} $\exists x\,\varphi(x)$.\\
\textbf{Take:} $t$ := [specific term].\\
\textbf{Verify:} $\varphi(t)$ holds. \quad [proof that $t$ works]\\
\textbf{Therefore:} $\exists x\,\varphi(x)$. \quad (by EG)
\end{quote}

The entire burden of an existential proof is finding the right $t$. Once $t$
is identified and verified, the logic is trivial: one application of EG.

\begin{remark}[Where the difficulty lies]
In existential proofs, the logical structure is simple but the mathematical
insight is often deep. Finding the witness requires understanding the problem;
verifying the witness is usually straightforward. This is why analysis proofs
often say ``take $\delta = \min(\varepsilon/2, 1)$'' without explaining where
that choice came from --- the work happened before the formal proof was written.
\end{remark}

% ---------------------------------------------------------
% Using a universal
% ---------------------------------------------------------

\section*{Strategy 3: Using a Universal $\forall x\,\varphi(x)$}

If a hypothesis or established fact has the form $\forall x\,\varphi(x)$, one
can instantiate it at \emph{any} specific term $t$ to obtain $\varphi(t)$, by
Universal Instantiation (UI). There are no side conditions beyond $t$ being
free for $x$ in $\varphi$.

\begin{quote}
\textbf{Available:} $\forall x\,\varphi(x)$.\\
\textbf{Choose any term $t$.}\\
\textbf{Obtain:} $\varphi(t)$. \quad (by UI)
\end{quote}

UI can be applied repeatedly to the same universal statement with different
choices of $t$.

% ---------------------------------------------------------
% Using an existential
% ---------------------------------------------------------

\section*{Strategy 4: Using an Existential $\exists x\,\varphi(x)$}

If a hypothesis has the form $\exists x\,\varphi(x)$, introduce a \emph{fresh
constant} $c$ to name the (unknown, unspecified) witness, obtaining $\varphi(c)$
by Existential Instantiation (EI). The constant $c$ must be new: it must not
appear anywhere else in the proof.

\begin{quote}
\textbf{Available:} $\exists x\,\varphi(x)$.\\
\textbf{Let $c$ be a witness.} (fresh constant, not used elsewhere)\\
\textbf{Obtain:} $\varphi(c)$. \quad (by EI)
\end{quote}

The freshness requirement is not a technicality. Using a constant that already
appears in the proof --- say, a constant $a$ introduced under a different
assumption --- would illegally identify the witness with that specific element.
The witness $c$ is only known to satisfy $\varphi(c)$; nothing else about it
is known, and it cannot be equated with anything specific.

\begin{remark}[The scope of a witness constant]
Once $c$ is introduced by EI, it names the witness for the duration of the
argument. Conclusions derived using $c$ depend on the existence of such a
witness; they are valid (they do not depend on which specific element $c$
actually names). However, $c$ cannot be exported outside the scope in which
the existential was instantiated, and no universal generalization may be
applied to $c$.
\end{remark}

% ---------------------------------------------------------
% The critical asymmetry: forall vs exists goals
% ---------------------------------------------------------

\section*{The Central Asymmetry: Universal vs.\ Existential Goals}

The following comparison is the most important thing to understand about
quantifier proof structure. Universal and existential goals require fundamentally
different strategies, and confusing them is one of the most common serious errors
in mathematical writing.

\begin{tcolorbox}[colback=gray!6, colframe=gray!40, arc=2pt,
  left=6pt, right=6pt, top=4pt, bottom=4pt,
  title={\small\textbf{Universal vs.\ Existential Goals: The Fundamental Difference}},
  fonttitle=\small\bfseries]
\small
\begin{tabular}{l l l}
\toprule
 & \textbf{$\forall x\,\varphi(x)$} & \textbf{$\exists x\,\varphi(x)$} \\
\midrule
\textbf{To prove:} & Fix \emph{arbitrary} $x$; prove $\varphi(x)$ & \emph{Choose} specific $t$; verify $\varphi(t)$ \\
\textbf{Who chooses $x$?} & The adversary (you cannot choose) & You (the prover) \\
\textbf{Main rule:} & UG & EG \\
\textbf{Burden:} & Works for every $x$ without assumptions & A single witness suffices \\
\textbf{Error:} & Assuming something special about $x$ & Failing to exhibit a concrete $t$ \\
\bottomrule
\end{tabular}
\end{tcolorbox}

\begin{remark}[Who controls the variable]
For a universal goal, the variable $x$ is chosen by the ``opponent'' of the
proof: the proof must work no matter what $x$ is. The prover has no control
over $x$; hence the term ``arbitrary.'' For an existential goal, the prover
chooses the witness $t$; the power is entirely on the prover's side, but so is
the obligation to produce something concrete and verifiable.
\end{remark}

% ---------------------------------------------------------
% Nested quantifiers: the analysis prototype
% ---------------------------------------------------------

\section*{Nested Quantifiers: The $\forall\exists$ Pattern in Analysis}

Every $\varepsilon$-$\delta$ proof in analysis is a proof of a statement with
the pattern $\forall\varepsilon\,\exists\delta\,\cdots$. Understanding the
proof structure of such statements is the direct payoff of this section.

\begin{tcolorbox}[colback=propbox, colframe=propborder, arc=2pt,
  left=6pt, right=6pt, top=4pt, bottom=4pt,
  title={\small\textbf{Proof Skeleton: $\forall x\,\exists y\,\varphi(x,y)$}},
  fonttitle=\small\bfseries]
\begin{enumerate}[label=\textbf{Step \arabic*.}]
\item \textbf{Fix arbitrary $x$.} \quad (Opening move for the outer $\forall$; $x$ is
  unspecified. This is Strategy 1.)
\item \textbf{Define $y$ in terms of $x$.} \quad (The witness for $\exists y$ may, and
  typically does, depend on $x$. This is Strategy 2 — choosing a witness.)
\item \textbf{Verify $\varphi(x,y)$} \quad using the specific $y$ from Step 2 and
  the arbitrary $x$ from Step 1.
\item \textbf{Conclude:} $\exists y\,\varphi(x,y)$. \quad (by EG)
\item \textbf{Conclude:} $\forall x\,\exists y\,\varphi(x,y)$. \quad (by UG, since $x$
  was arbitrary in Step 1)
\end{enumerate}
The key insight: $y$ is chosen \emph{after} $x$ is fixed, so $y$ is allowed
to depend on $x$.
\end{tcolorbox}

\begin{tcolorbox}[colback=propbox, colframe=propborder, arc=2pt,
  left=6pt, right=6pt, top=4pt, bottom=4pt,
  title={\small\textbf{Proof Skeleton: $\exists y\,\forall x\,\varphi(x,y)$}},
  fonttitle=\small\bfseries]
\begin{enumerate}[label=\textbf{Step \arabic*.}]
\item \textbf{Define $y$.} \quad (The witness for $\exists y$ must be chosen \emph{before}
  $x$ is introduced. This $y$ must work for \emph{every} $x$.)
\item \textbf{Fix arbitrary $x$.} \quad (Now $x$ is introduced.)
\item \textbf{Verify $\varphi(x,y)$} \quad for the specific $y$ from Step 1 and
  arbitrary $x$ from Step 2.
\item \textbf{Conclude:} $\forall x\,\varphi(x,y)$. \quad (by UG)
\item \textbf{Conclude:} $\exists y\,\forall x\,\varphi(x,y)$. \quad (by EG)
\end{enumerate}
The key contrast: $y$ is chosen \emph{before} $x$, so $y$ cannot depend on
$x$. This is why $\exists y\,\forall x$ is a much stronger claim.
\end{tcolorbox}

\begin{remark}[Why order determines strength]
In the $\forall\exists$ skeleton, the witness $y$ is constructed from $x$, so
a different $y$ is allowed for each $x$. In the $\exists\forall$ skeleton, a
single $y$ must work for all $x$ simultaneously. The first pattern appears
in continuity ($\delta$ may depend on $\varepsilon$ and $x$); the second in
uniform continuity ($\delta$ may depend on $\varepsilon$ but not on $x$).
The difference between continuity and uniform continuity is exactly the
difference between these two proof skeletons.
\end{remark}

% ---------------------------------------------------------
% Full worked example
% ---------------------------------------------------------

\section*{A Complete Worked Example}

The following proof traces every UI/UG/EI/EG step explicitly. In practice,
proofs are written more concisely, but the logical structure is always this one.

\begin{example}[A worked multi-quantifier proof]
\textbf{Claim:} If $\forall x\,\exists y\,P(x,y)$ and $\forall x\,\forall y\,
(P(x,y) \to Q(x,y))$, then $\forall x\,\exists y\,Q(x,y)$.

\medskip
\textbf{Proof:}

\begin{enumerate}
\item Assume $H_1 := \forall x\,\exists y\,P(x,y)$ and
  $H_2 := \forall x\,\forall y\,(P(x,y) \to Q(x,y))$.
  \quad (hypotheses)

\item Let $a$ be an arbitrary element of the domain.
  \quad (introduce arbitrary variable for $\forall x$ in the goal)

\item From $H_1$, instantiate at $a$: $\exists y\,P(a,y)$.
  \quad (by UI applied to $H_1$ with $t = a$)

\item Let $b$ be a witness: $P(a,b)$.
  \quad (by EI applied to Step 3; $b$ is fresh)

\item From $H_2$, instantiate at $a$: $\forall y\,(P(a,y) \to Q(a,y))$.
  \quad (by UI applied to $H_2$ with $t = a$)

\item From Step 5, instantiate at $b$: $P(a,b) \to Q(a,b)$.
  \quad (by UI applied to Step 5 with $t = b$)

\item From Step 4 ($P(a,b)$) and Step 6 ($P(a,b) \to Q(a,b)$): obtain $Q(a,b)$.
  \quad (by Modus Ponens)

\item From Step 7: $\exists y\,Q(a,y)$.
  \quad (by EG: $b$ witnesses the existential)

\item Since $a$ was arbitrary: $\forall x\,\exists y\,Q(x,y)$.
  \quad (by UG: $a$ was not constrained)
\end{enumerate}

\textbf{Annotation:} The proof follows the $\forall\exists$ skeleton exactly.
Step 2 fixes the arbitrary $x$. Steps 3--4 use the first hypothesis to obtain
a witness $b$ depending on $a$. Steps 5--7 apply the second hypothesis to get
the desired conclusion about $b$. Steps 8--9 wrap everything back up with
EG and UG.
\end{example}

\begin{remark}[Reading analysis proofs]
A proof that begins ``Let $\varepsilon > 0$'' is executing Step 2 of the
$\forall\exists$ skeleton with domain $\mathbb{R}_{>0}$. A sentence like
``take $\delta = \varepsilon/2$'' is executing Step 2 of the inner existential
proof: it is the witness for $\exists\delta$. Everything between ``let
$\varepsilon > 0$'' and ``take $\delta = \ldots$'' is scaffolding to find the
right witness. Everything after ``take $\delta = \ldots$'' is the verification
that the witness works. Once this structure is seen, every $\varepsilon$-$\delta$
proof becomes readable by template.
\end{remark}
